\section{Daten bearbeiten mit \texorpdfstring{\gls{SQL}}{SQL}}
Mit SQL kann man Daten nicht nur abfragen, sondern auch erstellen, verändern und löschen. Im folgenden Schauen wir uns dafür nötigen Befehle an.

\subsection{Eine neue Tabelle erstellen: \texorpdfstring{\lstinline|CREATE TABLE|}{CREATE TABLE}}
\label{subsec:create-table}
Eine neue Tabelle kann innerhalb einer Datenbank mit dem Befehl \lstinline|CREATE TABLE| angelegt werden:

\begin{lstlisting}
CREATE TABLE tabellenname(
  kolonne1 eigenschaften1,
  kolonne2 eigenschaften2,
  ...
)
\end{lstlisting}

\begin{myexample}
\label{ex:sql-create}
Die Tabelle \texttt{users} wurde mit folgendem Befehl erstellt:

\begin{lstlisting}
CREATE TABLE users (
  `id` int(10) UNSIGNED NOT NULL AUTO_INCREMENT PRIMARY KEY,
  `username` varchar(191) NOT NULL UNIQUE,
  `email` varchar(191) NOT NULL UNIQUE,
  `password` varchar(191) NOT NULL,
  `name` varchar(191) NOT NULL,
  `bio` varchar(191) DEFAULT NULL,
  `gender` enum('male','female') DEFAULT NULL,
  `birthday` datetime DEFAULT NULL,
  `city` varchar(191) DEFAULT NULL,
  `country` varchar(191) DEFAULT NULL,
  `centimeters` int(11) DEFAULT NULL,
  `avatar` varchar(191) NOT NULL DEFAULT 'avatar.png',
  `role` enum('user','dba','teacher','admin') NOT NULL DEFAULT 'user',
  `is_active` tinyint(1) NOT NULL DEFAULT 0,
  `remember_token` varchar(100) DEFAULT NULL,
  `created_at` timestamp NULL DEFAULT NULL,
  `updated_at` timestamp NULL DEFAULT NULL
)
\end{lstlisting}
\end{myexample}

Nach dem Kolonnennamen folgt zuerst der \textbf{Kolonnentyp}, d.h., die Angabe, welche Art von Daten in der Kolonne gespeichert werden. Einige der häufigsten Datentypen sind:
\begin{itemize}
  \item \lstinline|INT(laenge)| bezeichnet eine Ganzzahl, für die \texttt{laenge} bytes zur Verfügung stehen. Beispielsweise speichert die Kolonne \texttt{centimeters} eine Zahl von 0 bis $2^{11}$, d.h. 2048.
  \item \lstinline|TINYINT(1)| speichert den Wert 0 oder 1, wobei 0 typischerweise ``falsch'' und 1 ``wahr'' bedeuten.
  \item \lstinline|DATETIME| ist ein Datum mit einer Uhrzeit
  \item \lstinline|VARCHAR(laenge)| bezeichnet eine Zeichenkette, d.h. ein Text, der mit maximal \texttt{laenge} vielen Bit kodiert wird.
  \item \lstinline|TEXT| bezeichnet eine Zeichenkette mit maximaler Länge von 65'535 bytes.
  \item \lstinline|ENUM('wert1', 'wert2',...)| bezeichnet eine Auflistung (engl. \textit{enumeration}) möglicher Werte. Beispielsweise können nur ``male'' und ``female'' im Feld \texttt{gender} eingetragen werden.
\end{itemize}

Zudem können folgende Eingeschaften für Kolonnen angegeben werden:

\begin{itemize}
  \item \lstinline|NOT NULL|: Feld darf nicht leer sein. Falls kein \lstinline|DEFAULT|-Wert angegeben wird, muss, der Wert für diese Kolonne beim Erstellen eines neuen Benutzerkontos immer mitgeliefert werden.
  \item \lstinline|DEFAULT wert|: Standard-Wert, falls nichts anderes angegeben wird. Das Attribut \texttt{avatar} nimmt beispielsweise den Wert \texttt{avatar.png} an, falls kein anderes Bild hochgeladen wird.
  \item \lstinline|AUTO_INCREMENT|: Zahl, die bei jedem neuen Eintrag automatisch immer um 1 grösser wird
  \item \lstinline|PRIMARY KEY|: Primärschlüssel (s. \cref{table:table-cia}).
  \item \lstinline|FOREIGN KEY (kolonneID) REFERENCES tabelle2(kolonneID2)|: Fremdschlüssel (s. \cref{table:table-cia}).
  \item \lstinline|UNIQUE|: Darf keine doppelten Werte enthalten
  \item \lstinline|UNSIGNED|: Ohne Vorzeichen (+-), also nur positive Zahlen
\end{itemize}

\begin{mycaution}
	Je nach SQL-Version muss ein Fremdschlüssel erst \textit{nach dem Erstellen einer Tabelle} als Fremdschlüssel deklariert werden. Dies kann mit dem Befehl \lstinline|ALTER TABLE| gemacht werden. 
	\begin{myexample}
		In folgendem Beispiel wird in Tabelle \lstinline|table2| das Attribut \lstinline|fID| nachträglich als Fremdschlüssel deklariert, indem es auf das Attribut \lstinline|ID| bei \lstinline|table1| zeigt.
		\begin{lstlisting}
			CREATE TABLE table1 (
				ID int(10) UNSIGNED NOT NULL AUTO_INCREMENT PRIMARY KEY,
				...
			);
			
			CREATE TABLE table2 (
				`ID` int(10) UNSIGNED NOT NULL AUTO_INCREMENT PRIMARY KEY,
				`fID` int(10) UNSIGNED NOT NULL
				...
			);
			
			ALTER TABLE table2
			ADD FOREIGN KEY (fID) REFERENCES table1(ID)
		\end{lstlisting}
	\end{myexample}
\end{mycaution}

	\begin{mycaution}
Auf InstaHub kann immer nur ein Befehl pro Mal ausgeführt werden, das ``;''-Zeichen funktioniert also leider nicht. Stattdessen müssen Sie in InstaHub jeden Befehl einzeln eingeben.
	\end{mycaution}

\begin{myexercise}
	Schreiben Sie die SQL-Befehle, um \cref{table:angestellte2} und \cref{table:abteilungen} zu erstellen. Sie müssen dabei lediglich eine leere Tabelle mit den korrekten Kolonnen-Namen und Kolonnen-Typen (Zahl, Text, etc.) erstellen, Daten müssen Sie noch keine einfügen.
	\begin{myanswer}
		\begin{lstlisting}
			CREATE TABLE abteilungen (
				aID int(10) UNSIGNED NOT NULL AUTO_INCREMENT PRIMARY KEY,
				abteilungsname varchar(191) NOT NULL UNIQUE
			)
			\end{lstlisting}
			\begin{lstlisting}
			CREATE TABLE angestellte (
				`mID` int(10) UNSIGNED NOT NULL AUTO_INCREMENT PRIMARY KEY,
				`name` varchar(191) NOT NULL UNIQUE,
				`aID` int(10) UNSIGNED NOT NULL
			)
			\end{lstlisting}
		\begin{lstlisting}
			ALTER TABLE angestellte
			ADD FOREIGN KEY (aID) REFERENCES abteilungen(aID)
		\end{lstlisting}
	\end{myanswer}
\end{myexercise}

\subsection{Tabellen verändern: \texorpdfstring{\lstinline|ALTER|}{ALTER}}
Mit dem Befehl \lstinline|ALTER tabellenname| können Tabellen und deren Attribute nachträglich, d.h. nach der Erstellung der Tabelle, verändert werden.

\begin{myexample}
Um eine weitere Kolonne hinzuzufügen, kann folgender Befehl verwendet werden:
\begin{lstlisting}
ALTER TABLE tabellenname
ADD neue_kolonne eigenschaften
\end{lstlisting}
\end{myexample}

\begin{myexample}
Um eine Kolonne umzubenennen, kann folgender Befehl verwendet werden:
\begin{lstlisting}
ALTER TABLE tabellenname
RENAME COLUMN alter_name TO neuer_name;
\end{lstlisting}
\end{myexample}
\begin{mycaution}
Beim Umbenennen einer Kolonne können Datenbankschemen (z.B. Fremschlüssel-Referenzen) zerstört werden, z.B. falls ein Primärschlüssel umbenannt wird, welcher von einer anderen Tabelle per Fremschlüssel referenziert wird (s. \cref{subsec:create-table}).
\end{mycaution}

\subsection{Tabellen löschen: \texorpdfstring{\lstinline|DROP TABLE|}{DROP TABLE}}
Mit dem Befehl \lstinline|DROP TABLE tabellenname| kann eine Tabelle komplett löschen.

\begin{mycaution}
  Dieser Schritt ist irreversibel. Falls Sie eine Tabelle irrtümlich löschen, geben Sie mir Bescheid, damit ich Ihre Datenbank zurücksetzen kann.
\end{mycaution}

\begin{myexercise}
  Erstellen Sie auf InstaHub eine Tabelle \texttt{users2} mit denselben Eigenschaften wie \texttt{users} (s. Beispiel \ref{ex:sql-create})
  \begin{myanswer}
    \begin{lstlisting}
CREATE TABLE users2 (
  `id` int(10) UNSIGNED NOT NULL AUTO_INCREMENT PRIMARY KEY,
  `username` varchar(191) NOT NULL UNIQUE,
  `email` varchar(191) NOT NULL UNIQUE,
  `password` varchar(191) NOT NULL,
  `name` varchar(191) NOT NULL,
  `bio` varchar(191) DEFAULT NULL,
  `gender` enum('male','female') DEFAULT NULL,
  `birthday` datetime DEFAULT NULL,
  `city` varchar(191) DEFAULT NULL,
  `country` varchar(191) DEFAULT NULL,
  `centimeters` int(11) DEFAULT NULL,
  `avatar` varchar(191) NOT NULL DEFAULT 'avatar.png',
  `role` enum('user','dba','teacher','admin') NOT NULL DEFAULT 'user',
  `is_active` tinyint(1) NOT NULL DEFAULT 0,
  `remember_token` varchar(100) DEFAULT NULL,
  `created_at` timestamp NULL DEFAULT NULL,
  `updated_at` timestamp NULL DEFAULT NULL
)
\end{lstlisting}
  \end{myanswer}
\end{myexercise}

\begin{myexercise}
  Löschen Sie danach die Tabelle \texttt{users2} wieder.
  \begin{myanswer}
    \begin{lstlisting}
DROP TABLE users2
\end{lstlisting}
  \end{myanswer}
\end{myexercise}


\subsection{Daten einfügen: \texorpdfstring{\lstinline|INSERT INTO|}{INSERT INTO}}

Mit \lstinline|INSERT INTO| können neue Zeilen in eine bestehende Tabelle eingefügt werden.

Die Syntax von \lstinline|INSERT INTO| sieht wie folgt aus:

\begin{lstlisting}
INSERT INTO tabelle_name (kolonne1, kolonne2, kolonne3, ...)
VALUES (wert1, wert2, wert3, ...)
\end{lstlisting}

\begin{myexample}
\label{ex:sql-insert}
Folgender Code fügt eine neue Person in die Tabelle \texttt{users} ein:
\begin{lstlisting}
INSERT INTO users (username, email, password, name, bio, gender, birthday, city, country, centimeters, avatar, role, is_active, remember_token, created_at, updated_at) 
VALUES ('guenther37', 'guenther@instahub.test', '12345', 'Günther Müller', 'Günther mag Kartoffelsalat.', 'male', '2006-06-06 00:00:00', 'Leipzig', 'Deutschland', '173', 'avatar.png', 'user', '0', NULL, now(), now())
\end{lstlisting}
Wie Sie sehen, wird zuerst die Tabelle (\texttt{users}) aufgelistet, danach die Kolonnen und schliesslich die Werte.
\end{myexample}

\begin{myexercise}
\label{aufg:sql-insert}
  Für welche Kolonnen müssen für einen neuen Benutzer immer die Werte mitgegeben werden? S. die Erklärungen zu \lstinline|NOT NULL| und \lstinline|DEFAULT| in \cref{subsec:create-table}.
  \begin{myanswer}
    \begin{verbatim}
username, email, password, name, role, is_active
\end{verbatim}
  \end{myanswer}
\end{myexercise}
\begin{mysubmission}\label{aufg:a-sql-insert}
  Erstellen Sie einen neuen Eintrag in der Tabelle \texttt{users} mit folgenden Eigenschaften:
  \begin{table}[H]
    \centering
    \begin{tabular}{c|c}
      \textbf{Kolonne}    & \textbf{Wert}                       \\\midrule
      \texttt{id}         & \texttt{300}                        \\
      \texttt{username}   & \texttt{testuser}                   \\
      \texttt{email}      & \texttt{test@test.com}              \\
      \texttt{password}   & \texttt{test123}                    \\
      \texttt{name}       & \texttt{Test-Vorname Test-Nachname} \\
      \texttt{role}       & \texttt{user}                       \\
      \texttt{is\_active} & \texttt{1}                          \\
    \end{tabular}
  \end{table}
  Überprüfen Sie danach, ob der Eintrag korrekt erstellt worden ist, indem Sie die Daten des users \texttt{testuser} abfragen (alle Attribute).
  \begin{myanswer}
    \begin{lstlisting}
INSERT INTO users (id, username, email, password, 
				   name, role, is_active) 
VALUES (300, 'testuser', 'test@test.com', 'test123', 'Test-Vorname Test-Nachname', 'user', '1')
\end{lstlisting}
Das Resultat kann mit einer zweiten Abfrage überprüft werden:
\begin{lstlisting}
SELECT * FROM users WHERE username="testuser"
\end{lstlisting}
  \end{myanswer}
\end{mysubmission}

\subsection{Daten verändern: \texorpdfstring{\lstinline|UPDATE|}{UPDATE}}

Mit \lstinline|UPDATE| können Werte einer Tabelle aktualisiert werden. Die Syntax von \lstinline|UPDATE| sieht wie folgt aus:

\begin{lstlisting}
UPDATE tabelle_name
SET kolonne1 = wert1, kolonne2=wert2, ...
WHERE bedingung(en)
\end{lstlisting}

\begin{myexercise}\label{aufgabe:evalPoly}

  Ersetzen Sie die Stadt ``Berlin'' überall durch ``Bern''
  \begin{myanswer}
    \begin{lstlisting}
UPDATE users
SET city="Bern"
WHERE city="Berlin"
\end{lstlisting}
  \end{myanswer}
\end{myexercise}

\begin{mysubmission}
  Setzen Sie die Körpergrössen aller männlichen Mitglieder auf 190. Überprüfen Sie danach die Richtigkeit Ihres Befehls, die Kolonnen für die Körpergrösse, das Geschlecht und den Name aller männlichen Mitglieder abfragen.
  \begin{myanswer}
    \begin{lstlisting}
UPDATE users
SET centimeters=190
WHERE gender="male"
\end{lstlisting}
Um zu überprüfen, ob der Befehl funktioniert hat, können Sie danach einfach die Tabelle nochmals anzeigen:
\begin{lstlisting}
SELECT centimeters, gender, name
FROM users
WHERE gender="male"
\end{lstlisting}

  \end{myanswer}
\end{mysubmission}

\subsection{Einträge löschen: \texorpdfstring{\lstinline|DELETE FROM|}{DELETE FROM}}
Mit dem Befehl \lstinline|DELETE FROM| können Einträge aus einer Tabelle gelöscht werden.
\begin{lstlisting}
DELETE FROM tabellenname
WHERE bedingungen
\end{lstlisting}

\begin{myexample}
Folgender Code löscht den Eintrag für user \texttt{guenther37}:
\begin{lstlisting}
DELETE FROM users
WHERE username="guenther37"
\end{lstlisting}
\end{myexample}

\begin{myexercise}
  Löschen Sie den Eintrag, den Sie in \cref{aufg:a-sql-insert} erstellt haben. Überprüfen Sie, ob der Eintrag verschwunden ist, indem Sie danach \lstinline|SELECT * FROM users| ausführen.
  \begin{myanswer}
    \begin{lstlisting}
DELETE FROM users
WHERE username="testuser"
\end{lstlisting}
  \end{myanswer}
\end{myexercise}

\section{Weiterführende Links und Übungen}
Folgende Links dienen der Vertiefung in das Thema \gls{SQL} (und der Prüfungsvorbereitung).

Allgemein zu Datenbanken:
\begin{itemize}
  \item \url{https://oinf.ch/kurs/vernetzung-und-systeme/datenbanken/}
  \item \url{https://wi-wissen.github.io/instahub-doc-de/#/exercices}
\end{itemize}
Zu \gls{SQL}:
\begin{itemize}
  \item \url{https://www.w3schools.com/sql/default.asp} (auf Englisch)
  \item \url{https://sql-island.informatik.uni-kl.de}
  \item \url{https://sql-tutorial.de/home/lektionen.php?lektion=1}. Hier wird unter anderem folgende Datenbank verwendet:
        \begin{figure}[H]
          \centering
          \begin{tikzpicture}[entity/.style={ellipse, draw=black, minimum width=2cm}, key/.style={rectangle, draw=black, minimum width=2cm, minimum height=1cm}]

            % Entity "laender"
            \node[key] (laender) at (0,0) {cia};

            % Attribute "Name" with underlining to denote a key
            \node[entity, left=of laender] (name) {\underline{name}};
            \draw (laender) -- (name);

            % Attribute "Region"
            \node[entity, above left=of laender] (region) {region};
            \draw (laender) -- (region);

            % Attribute "Fläche"
            \node[entity, above=of laender] (flaeche) {flaeche};
            \draw (laender) -- (flaeche);

            % Attribute "Einwohner"
            \node[entity, above right=of laender] (einwohner) {einwohner};
            \draw (laender) -- (einwohner);

            % Attribute "BIP"
            \node[entity, right=of laender] (bip) {bip};
            \draw (laender) -- (bip);

          \end{tikzpicture}
        \end{figure}
\end{itemize}
